\scanpage{121}
\mspage{5}
% archive sheet number: 120
\diagram{\begin{tikzcd}[ampersand replacement=\&,row sep=large,column sep=small]
\& X_{\unsure{1|4}} \arrow[dl,Rightarrow] \arrow[dd,"p_1",near start] \& \& X_3 \arrow[ll] \arrow[dl,dash] \arrow[dd,dotted,"p_3"] \\
X_0 \arrow[dd,"p_0"'] \& \& X_2 \arrow[ll] \arrow[dd,"p_2"'] \& \\
\& Y_1 \arrow[dl] \& \& Y_3 \arrow[ll] \arrow[dl] \\
Y_0 \& \& Y_2 \arrow[ll] \&
\end{tikzcd}}

Alors :

a) il existe un hom. $\bar p:X\to Y$ induisant $p$, i.e. il existe une flèche $p_3$ de façon à compléter ce diagramme cubique commutatif (en respectant donc les contraintes de commutativité pour les deux carrés adjacents à l'arête $p_3$).

b) Si de plus $p_0,p_1,p_2$ sont des iso, il en est de même de $p_3$.

Prenant le cas où $p_0,p_1,p_2$ sont des identités, on trouve :

\uline{Corollaire 1} dans les conditions de $\mathrm{Cart}(4)$, si on a \uline{deux} carrés $h$-cartésiens \inserted{$X,X'$} qui complètent
\diagram{\begin{tikzcd}[ampersand replacement=\&]
\& \& X_1 \arrow[dl] \\
\deleted{\unsure{X_0}} \& X_0 \& \\
\& \& X_2 \arrow[ul]
\end{tikzcd}}
, il existe un hom. de $\unsure{\alpha|X}$ dans $\alpha'$ qui induise l'identité sur les sommets d'indices $0,1,2$, et un tel hom. est tjrs un isomorphisme.

Mais cet isom. n'est pas forcément unique, car \unsure{d'autres|d'entrée|d'autre} \unsure{dans|de jeu|part} un autom. d'un carré \deleted{\unsure{les}} \inserted{$h$-}cartésien peut induire l'identité sur les sommets $X_0,X_1,X_2$, sans être l'identité.
